% Generated standalone source. Edit cv/profile.json or cv/template.tex and regenerate.
\documentclass[11pt,a4paper]{article}
\usepackage[margin=16mm]{geometry}
\usepackage{fontspec}
\setmainfont{lmroman10-regular.otf}[BoldFont=lmroman10-bold.otf,ItalicFont=lmroman10-italic.otf,BoldItalicFont=lmroman10-bolditalic.otf,Ligatures=NoCommon]
\usepackage{enumitem,titlesec}
\usepackage[hidelinks,unicode]{hyperref}
\hypersetup{pdftitle={Javier Gomez Mijangos - General Research},pdfauthor={Javier Gomez Mijangos},pdfsubject={Undergraduate internship application}}
\pagestyle{empty}
\hyphenpenalty=10000
\exhyphenpenalty=10000
\setlength{\parindent}{0pt}
\setlength{\parskip}{0pt}
\setlength{\emergencystretch}{1em}
\titleformat{\section}{\large\bfseries}{}{0pt}{}[\titlerule]
\titlespacing*{\section}{0pt}{9pt}{5pt}
\setlist[itemize]{leftmargin=1.2em,itemsep=2pt,topsep=3pt,parsep=0pt,partopsep=0pt}
\newcommand{\entry}[2]{\textbf{#1}\enspace|\enspace{\small #2}\par}
\begin{document}
\raggedright
\begin{center}
{\LARGE\bfseries Javier Gomez Mijangos}\\[4pt]
Mathematics \& Computer Science | Undergraduate Research Internships\\[4pt]
\href{mailto:javier.gomezmijan@gmail.com}{javier.gomezmijan@gmail.com}\\[2pt]
\href{https://github.com/Javpato}{github.com/Javpato} | \href{https://www.linkedin.com/in/javier-gomez-a01151283/}{linkedin.com/in/javier-gomez-a01151283} | \href{https://javpato.github.io/}{javpato.github.io}
\end{center}
\section{Education}
\entry{Université Paris-Saclay}{Sep 2025 - Jun 2028 (expected)}
Double Bachelor's Degree in Mathematics and Computer Science\\[2pt]
Average: \textbf{16.72/20}.
\section{Selected Projects}
\entry{\href{https://github.com/Javpato/RayTracingProject}{C++ Ray Tracer}}{2025}
{\small Independent tutorial-based project}\par
{\small C++17, CMake, 3D geometry}\par
\begin{itemize}
\item Implemented ray-sphere intersections and recursive light scattering in C++17 while following Ray Tracing in One Weekend.
\item Translated vector geometry, reflection and refraction into material models, including Schlick approximation for dielectric reflectance.
\end{itemize}
\entry{\href{https://github.com/Javpato/Data-Science-Spoon-knives}{Utensil Image Classification}}{2026}
{\small University project}\par
{\small Python, NumPy, pandas, scikit-learn}\par
\begin{itemize}
\item Built an image-processing pipeline to isolate utensils and extract six geometric features, including elongation, roundness and template-matching scores.
\item Compared KNN, SVM, tree-based and neural-network classifiers using balanced accuracy and confusion matrices on a small knife/spoon dataset.
\item Visualized class-specific feature distributions to examine how geometric descriptors separate knives from spoons.
\end{itemize}
\entry{\href{https://github.com/Javpato/diffusion-model}{Diffusion and Image Denoising}}{2026}
{\small Independent course-based project}\par
{\small Python, PyTorch, Hugging Face diffusers}\par
\begin{itemize}
\item Implemented denoising training loops on MNIST/Fashion-MNIST while studying an MIT course, using PyTorch and Hugging Face U-Net components.
\item Trained noise prediction with sampled diffusion timesteps, mean-squared error and Adam; generated images through iterative DDPM denoising.
\item Recorded training-loss curves and inspected generated samples to examine how the model denoises images.
\end{itemize}
\section{Technical Skills}
Scientific programming: Python, C++, NumPy, pandas, PyTorch, scikit-learn.\\[2pt]
Image processing, numerical geometry, model diagnostics; Git, Linux.
\section{Award}
\textbf{1st Place, Build Sprint (BuildNXT)} \enspace|\enspace Nov 2026\par
Battle Tower: programmer on a team of two, implementing all application features.
\section{Languages}
Spanish (native), English (fluent), French (fluent), Mandarin Chinese (conversational).
\end{document}
